% !TeX program = lualatex
% halo:begin
% {
%   "version": 1,
%   "publish": true,
%   "course": "mathematical-analysis-i",
%   "section": "6.3",
%   "title": "数学分析（一）6.3：泰勒公式",
%   "categories": ["study-notes"],
%   "tags": ["mathematics"],
%   "excerpt": "从导数确定泰勒多项式，用佩亚诺余项计算极限与高阶导数，用拉格朗日余项估计误差并证明 e 的无理性。"
% }
% halo:end
% 来源：IMG_3594 右页、IMG_3595；华东师范大学《数学分析》第5版上册，第六章§3，书页125–132（PDF 136–143）。
\RequirePackage{iftex}
\RequireLuaTeX
\documentclass[UTF8,fontset=fandol,a4paper,12pt]{ctexart}
\usepackage[margin=2.5cm]{geometry}
\usepackage{amsmath,amssymb,amsthm}
\usepackage{fancyhdr}
\usepackage[hidelinks]{hyperref}
\theoremstyle{definition}
\newtheorem{theorem}{定理}[section]
\newtheorem{proposition}[theorem]{命题}
\newtheorem{example}[theorem]{例题}
\renewcommand{\proofname}{证明}
\setlength{\headheight}{16pt}
\pagestyle{fancy}\fancyhf{}
\fancyhead[L]{\small 数学分析（一）}
\fancyhead[R]{\small 6.3\quad 泰勒公式}
\fancyfoot[C]{\thepage}\renewcommand{\headrulewidth}{0.3pt}
\raggedbottom
\title{泰勒公式}
\author{Yuchen Zhang}
\date{整理日期：2026年11月26日}
\makeatletter
\renewcommand{\maketitle}{\begin{center}{\LARGE\bfseries\@title\par}\vspace{0.6em}\@author\par\vspace{0.2em}{\small\@date\par}\end{center}}
\makeatother
\begin{document}
\maketitle

\section{在展开点配齐各阶导数}
导数给出的是最简单的一阶局部近似：若 $f$ 在 $x_0$ 处可导，令 $h=x-x_0$，则有 $f(x_0+h)=f(x_0)+f'(x_0)h+o(h)$。从几何上看，这条切线在 $x_0$ 处只与曲线对齐了高度和斜率；要想在展开点附近贴合得更紧，自然的做法是提高多项式次数，让它在 $x_0$ 处的各阶导数依次与 $f$ 吻合。

为此，设 $n$ 次多项式 $T_n(x)=\sum_{k=0}^{n}a_k(x-x_0)^k$。对 $T_n$ 逐次求导并代入 $x=x_0$，次数低于 $k$ 的项在求导中已经消去，次数高于 $k$ 的项仍含有因子 $x-x_0$ 而在 $x_0$ 处为零，因此恰好有 $T_n^{(k)}(x_0)=k!\,a_k$。令它等于 $f^{(k)}(x_0)$，各项系数便被唯一确定：
\[
 a_k=\frac{f^{(k)}(x_0)}{k!},\qquad
 T_n(x)=\sum_{k=0}^{n}\frac{f^{(k)}(x_0)}{k!}(x-x_0)^k.
\]
这里约定 $f^{(0)}=f$。这个唯一与 $f$ 在 $x_0$ 处具有相同 $0$ 至 $n$ 阶导数的多项式称为 $f$ 在 $x_0$ 处的\textbf{泰勒多项式}，系数 $a_k$ 称为泰勒系数，而差值 $R_n(x)=f(x)-T_n(x)$ 称为余项。

接下来的定理说明，这种按导数配齐的多项式确实能在展开点附近给出高阶局部近似。

\begin{theorem}[带佩亚诺型余项的泰勒公式，教材定理 6.9]
设 $n\ge1$。若 $f$ 在 $x_0$ 处具有直到 $n$ 阶的导数，则当 $h\to0$ 时，
\[
 f(x_0+h)=\sum_{k=0}^{n}\frac{f^{(k)}(x_0)}{k!}h^k+o(h^n).
\]
\end{theorem}

这里的记号 $R_n(x_0+h)=o(h^n)$ 表示当 $h\to0$ 时有 $R_n(x_0+h)/h^n\to0$。例如 $h^2=o(h)$，因为 $h^2/h=h\to0$；相比之下，$0.001h$ 虽然同样趋于零，但它除以 $h$ 后恒等于 $0.001$，并不属于 $o(h)$。因此，仅知道误差趋于零还不足以判断它达到的阶数。

\begin{proof}
由泰勒系数的取法可知，余项满足 $R_n^{(k)}(x_0)=0$（$0\le k\le n$）。由于 $f^{(n)}(x_0)$ 存在，$R_n^{(n-1)}$ 在 $x_0$ 的某邻域内存在且在 $x_0$ 处可导，先由导数定义考察最高一阶的商：
\[
 \lim_{x\to x_0}\frac{R_n^{(n-1)}(x)}{n!(x-x_0)}
 =\frac{1}{n!}\lim_{x\to x_0}\frac{R_n^{(n-1)}(x)-R_n^{(n-1)}(x_0)}{x-x_0}
 =\frac{R_n^{(n)}(x_0)}{n!}=0.
\]
当 $n=1$ 时，上式已经直接给出结论。当 $n\ge2$ 时，对商式 $R_n(x)/(x-x_0)^n$ 逐次求导，前 $n-1$ 步所得均为 $0/0$ 型，且第 $n-1$ 次求导后恰好化为上式左端的极限。由洛必达法则自下而上逐层反推，即得 $\lim_{x\to x_0}R_n(x)/(x-x_0)^n=0$，亦即 $R_n(x_0+h)=o(h^n)$。
\end{proof}

需要指出，使误差达到 $o(h^n)$ 的不超过 $n$ 次的多项式是唯一的。事实上，若两个这样的多项式都满足误差为 $o(h^n)$，则它们之差必也是 $o(h^n)$。假如该差多项式非零，记其最低次非零项为 $c h^j$（$c\ne0$，$0\le j\le n$），两端同除以 $h^j$ 后令 $h\to0$，一端趋于 $c\ne0$，另一端却为 $o(h^{n-j})$ 而趋于零，矛盾。因此，只要已知 $f$ 在 $x_0$ 处 $n$ 阶可导，无论通过何种途径得到满足误差为 $o(h^n)$ 的多项式，都可以直接从中读出 $f$ 的泰勒系数。

\section{从近似多项式到麦克劳林公式}
\subsection{误差达到高阶，能否反推高阶可导？}
在 $n=1$ 的情形，若 $f(x)=a_0+a_1x+o(x)$ 且 $f(0)=a_0$，则由差商极限立即可得 $f'(0)=a_1$。那么，当误差阶数提高到 $o(x^n)$ 时，能否反过来依次推出函数在展开点具有二阶乃至更高阶导数？答案是否定的。

设 $n\ge2$，令狄利克雷函数 $D(x)$ 在有理数处取 $1$、在无理数处取 $0$，构造 $f(x)=x^{n+1}D(x)$。一方面，由于 $|f(x)/x^n|=|x|D(x)\le|x|\to0$，零多项式就已经使误差达到 $o(x^n)$，同时由 $f(x)/x=x^n D(x)\to0$ 可知 $f'(0)=0$。另一方面，对任意 $a\ne0$，分别沿有理数和无理数趋于 $a$ 时，$f(x)$ 的极限依次为 $a^{n+1}$ 与 $0$，二者不等，故 $f$ 在原点以外处处不连续。这意味着 $f'$ 在原点附近任何非零点处都没有定义，自然无法再讨论 $f''(0)$。由此可见，只有在另行确认函数的高阶可导性之后，才能把近似多项式的系数解释为各阶导数。

\subsection{六种常用的麦克劳林公式}
在泰勒公式中取展开点 $x_0=0$，便得到麦克劳林公式。以 $f(x)=\sin x$ 为例，其 $k$ 阶导数为 $f^{(k)}(x)=\sin(x+k\pi/2)$，代入 $x=0$ 后，偶数阶导数全为零，奇数阶导数正负交替。类似地计算各基本初等函数在原点处的导数，即可得到当 $x\to0$ 时的六种常用展开（其中 $m,n\ge1$）：
\begin{align*}
 e^x &= \sum_{k=0}^n\frac{x^k}{k!}+o(x^n),&
 \frac1{1-x}&=\sum_{k=0}^n x^k+o(x^n),\\
 \sin x &=\sum_{k=0}^{m-1}\frac{(-1)^k x^{2k+1}}{(2k+1)!}+o(x^{2m}),&
 \cos x &=\sum_{k=0}^{m}\frac{(-1)^k x^{2k}}{(2k)!}+o(x^{2m+1}),\\
 \ln(1+x)&=\sum_{k=1}^{n}(-1)^{k-1}\frac{x^k}{k}+o(x^n),&
 (1+x)^\alpha&=\sum_{k=0}^n\binom{\alpha}{k}x^k+o(x^n).
\end{align*}
上式中广义二项式系数定义为 $\binom{\alpha}{0}=1$，$\binom{\alpha}{k}=\alpha(\alpha-1)\cdots(\alpha-k+1)/k!$。以对数函数为例，其 $k$ 阶导数为 $(-1)^{k-1}(k-1)!/(1+x)^k$，在 $x=0$ 处取值并除以 $k!$ 便得到系数 $(-1)^{k-1}/k$。此外，正弦的 $2m$ 次项系数与余弦的 $2m+1$ 次项系数均为零，因此分别应用 $2m$ 阶与 $2m+1$ 阶泰勒公式，可以在不增加多项式项数的前提下，将余项的阶数各自提高一阶。

\section{代换、系数比较与极限}
有了基本展开式与近似多项式的唯一性，在求复合函数或隐式关系所确定函数的泰勒展开时，往往不必直接反复求导，而是通过变量代换与待定系数来简化运算。

\begin{example}
设 $f(x)=e^{-x^2/2}$，写出其麦克劳林公式，并求 $f^{(98)}(0)$ 与 $f^{(99)}(0)$。
\end{example}

\begin{proof}[解]
在 $e^u$ 的 $n$ 阶麦克劳林展开中令 $u=-x^2/2$，得
\[
 f(x)=\sum_{k=0}^{n}\frac{(-1)^k}{2^k k!}x^{2k}+o(x^{2n}),\qquad x\to0.
\]
这里利用了 $o(u^n)=o(x^{2n})$。因为 $f$ 在原点无穷次可导，由近似多项式的唯一性可知，上式正是 $f$ 的 $2n$ 阶麦克劳林公式。取 $n=49$，比较 $x^{98}$ 项的系数即得 $f^{(98)}(0)/98!=-1/(2^{49}\cdot 49!)$；又因 $f$ 是光滑偶函数，其奇数阶导函数均为奇函数，在原点处必为零。因此
\[
 f^{(98)}(0)=-\frac{98!}{2^{49}\,49!},\qquad f^{(99)}(0)=0.
\]
\end{proof}

在求高阶抵消型极限时，将分子展开到首项不再相互抵消的阶数，往往比多次使用洛必达法则更加清晰。

\begin{example}
求 $\displaystyle\lim_{x\to0}\frac{\cos x-e^{-x^2/2}}{x^4}$。
\end{example}

\begin{proof}[解]
结合上例的结果与余弦函数的麦克劳林展开，当 $x\to0$ 时有
\[
 e^{-x^2/2}=1-\frac{x^2}{2}+\frac{x^4}{8}+o(x^4),\qquad
 \cos x=1-\frac{x^2}{2}+\frac{x^4}{24}+o(x^4).
\]
两式相减后，常数项与二次项恰好抵消，首项为 $(1/24-1/8)x^4$，因此所求极限为 $1/24-1/8=-1/12$。
\end{proof}

若展开中心不在原点，可先作平移，将表达式改写为关于 $x-x_0$ 的基本形式再套用麦克劳林公式。

\begin{example}
写出 $\ln x$ 在 $x=2$ 处的 $n$ 阶泰勒公式。
\end{example}

\begin{proof}[解]
将 $\ln x$ 变形为 $\ln 2+\ln\bigl(1+(x-2)/2\bigr)$，在对数函数的麦克劳林展开中以 $(x-2)/2$ 替代自变量，即得
\[
 \ln x=\ln 2+\sum_{k=1}^n\frac{(-1)^{k-1}}{k\,2^k}(x-2)^k+o\bigl((x-2)^n\bigr),\qquad x\to2.
\]
特别地，当 $n\ge2$ 时，其一次项为 $(x-2)/2$，二次项为 $-(x-2)^2/8$。
\end{proof}

对于商函数或反函数，利用恒等式比较同次幂系数同样可以避开繁复的高阶求导。

\begin{example}[习题 6.3 第 1 题]
求 $\tan x$ 与 $\arctan x$ 在原点的五阶麦克劳林公式。
\end{example}

\begin{proof}[解]
这两个函数在原点附近均光滑且为奇函数，在原点处的一阶导数都等于 $1$，因而它们的展开式中不含偶次项，一次项系数均为 $1$。

先设 $\tan x=x+ax^3+bx^5+o(x^5)$，代入恒等式 $\cos x\tan x=\sin x$ 并展开至五次项，比较 $x^3$ 与 $x^5$ 的系数，得
\[
 a-\tfrac12=-\tfrac16,\qquad b-\tfrac a2+\tfrac1{24}=\tfrac1{120},
\]
解得 $a=1/3$、$b=2/15$，于是
\[
 \tan x=x+\tfrac13x^3+\tfrac{2}{15}x^5+o(x^5).
\]

再设 $u=\arctan x=x+cx^3+dx^5+o(x^5)$。由于当 $x\to0$ 时有 $u/x\to1$，代入 $\tan u$ 的展开式时满足 $o(u^5)=o(x^5)$。由 $\tan u=x$ 比较 $x^3$ 与 $x^5$ 的系数，得 $c+1/3=0$ 及 $d+c+2/15=0$，解得 $c=-1/3$、$d=1/5$，故
\[
 \arctan x=x-\tfrac13x^3+\tfrac15x^5+o(x^5).
\]
\end{proof}

\section{拉格朗日余项与误差估计}
佩亚诺余项只刻画 $x\to x_0$ 时误差的高阶趋零性质，适合求极限，却无法在具体某一点给出误差的数值界。例如用 $e^x$ 在 $0$ 处的展开近似计算 $e$ 时，需要在 $x=1$ 处估计余项的大小；这就不能只依赖展开点一处的导数，而必须借助 $x_0$ 与 $x$ 之间整个区间上的高阶导数。

\begin{theorem}[带拉格朗日型余项的泰勒公式，教材定理 6.10]
设 $n\ge0$，$x\ne x_0$。若 $f$ 在以 $x_0,x$ 为端点的闭区间上具有直到 $n$ 阶的连续导数，并在相应的开区间内具有 $n+1$ 阶导数，则严格介于 $x_0$ 与 $x$ 之间存在一点 $\xi$，使
\[
 R_n(x)=\frac{f^{(n+1)}(\xi)}{(n+1)!}(x-x_0)^{n+1}.
\]
\end{theorem}

\begin{proof}
固定 $x$，将展开中心视为变量 $t$，让 $t$ 从 $x_0$ 连续变化到 $x$，并构造两个辅助函数
\[
 F(t)=f(x)-\sum_{k=0}^n\frac{f^{(k)}(t)}{k!}(x-t)^k,\qquad G(t)=(x-t)^{n+1}.
\]
对 $F(t)$ 求导时，第 $k$ 项中对 $f^{(k)}(t)$ 求导所得的项，与第 $k+1$ 项中对 $(x-t)^{k+1}$ 求导所得的项恰好正负抵消，最终只余下末项的导数，即
\[
 F'(t)=-\frac{f^{(n+1)}(t)}{n!}(x-t)^n.
\]
又注意到 $F(x)=G(x)=0$ 且 $F(x_0)=R_n(x)$，在以 $x_0,x$ 为端点的闭区间上对 $F,G$ 应用柯西中值定理，存在严格介于 $x_0$ 与 $x$ 之间的点 $\xi$，使得
\[
 \frac{R_n(x)}{(x-x_0)^{n+1}}
 =\frac{F(x_0)-F(x)}{G(x_0)-G(x)}
 =\frac{F'(\xi)}{G'(\xi)}
 =\frac{f^{(n+1)}(\xi)}{(n+1)!}.
\]
移项即得拉格朗日型余项公式。
\end{proof}

若将中间点写为 $\xi=x_0+\theta(x-x_0)$（$0<\theta<1$），当 $n=0$ 时，上述公式便退化为拉格朗日中值公式 $f(x)-f(x_0)=f'(\xi)(x-x_0)$；当 $x_0=0$ 时，$\xi=\theta x$，即得带拉格朗日型余项的麦克劳林公式。在误差估计中，只要 $f^{(n+1)}$ 在 $x_0$ 与 $x$ 之间的开区间上满足 $|f^{(n+1)}|\le M$，便立刻得到定量的误差上界
\[
 |R_n(x)|\le\frac{M}{(n+1)!}|x-x_0|^{n+1}.
\]

\subsection{六种常用展开的拉格朗日余项}
沿用第 2 节的麦克劳林多项式，以下列出六种常见函数对应的拉格朗日型余项（其中 $n,m\ge1$，各式的 $\theta\in(0,1)$ 彼此无关）。对于末行的 $1/(1-x)$，教材给出的是 $|x|<1$，但就有限阶泰勒公式而言，只要从 $0$ 到 $x$ 的区间避开奇点 $1$，定理条件便已满足，因此适用范围可写为 $x<1$。
\begin{align*}
 e^x\colon\quad R_n(x)
   &=\frac{e^{\theta x}}{(n+1)!}x^{n+1}, && x\in\mathbb R,\\
 \sin x\colon\quad R_{2m}(x)
   &=\frac{(-1)^m\cos(\theta x)}{(2m+1)!}x^{2m+1}, && x\in\mathbb R,\\
 \cos x\colon\quad R_{2m+1}(x)
   &=\frac{(-1)^{m+1}\cos(\theta x)}{(2m+2)!}x^{2m+2}, && x\in\mathbb R,\\
 \ln(1+x)\colon\quad R_n(x)
   &=\frac{(-1)^n x^{n+1}}{(n+1)(1+\theta x)^{n+1}}, && x>-1,\\
 (1+x)^\alpha\colon\quad R_n(x)
   &=\binom{\alpha}{n+1}(1+\theta x)^{\alpha-n-1}x^{n+1}, && x>-1,\quad\alpha\in\mathbb R,\\
 \frac{1}{1-x}\colon\quad R_n(x)
   &=\frac{x^{n+1}}{(1-\theta x)^{n+2}}, && x<1.
\end{align*}

\section{用误差界选择阶数与允许范围}
在具体的数值与理论问题中，拉格朗日余项主要有两类用法：一是固定自变量取值，通过误差界确定展开阶数或导出数论性质；二是固定近似多项式的阶数，反推在给定精度要求下自变量的允许范围。

\begin{example}
计算 $e$ 的近似值，使绝对误差小于 $10^{-6}$，并证明 $e$ 为无理数。
\end{example}

\begin{proof}[解]
记部分和 $S_n=\sum_{k=0}^n 1/k!$。在 $e^x$ 的拉格朗日余项展开式中令 $x=1$，存在 $\theta\in(0,1)$ 使
\[
 e=S_n+\frac{e^\theta}{(n+1)!}.
\]
先估计 $e$ 的粗略上界：取 $n=2$，由 $e^\theta<e$ 得 $e<5/2+e/6$，解得 $e<3$。于是对任意 $n$ 均有 $0<e-S_n<3/(n+1)!$。令 $n=9$，则 $0<e-S_9<3/10!<10^{-6}$。计算 $S_9=98641/36288$ 可得
\[
 2.7182815<S_9<e<S_9+\frac{3}{10!}<2.7182824.
\]
因此取 $e\approx2.718282$，即使计入末位舍入误差，其绝对误差仍小于 $5\times10^{-7}$。

再证 $e$ 为无理数。反设 $e=p/q$，其中 $p,q$ 为正整数。取整数 $n\ge\max\{q,2\}$，则 $n!\,e$ 与 $n!\,S_n=\sum_{k=0}^n n!/k!$ 均为整数。然而根据上面的余项估计，
\[
 0<n!(e-S_n)=\frac{e^\theta}{n+1}<\frac{3}{n+1}\le1,
\]
这表明两个整数之差严格落在 $(0,1)$ 内，矛盾。故 $e$ 为无理数。
\end{proof}

\begin{example}
分别用 $x$ 和 $x-x^3/6$ 近似 $\sin x$，求使绝对误差不超过 $10^{-3}$ 的角度范围。
\end{example}

\begin{proof}[解]
用一次式 $x$ 近似 $\sin x$ 时，因二次项系数为零，可在正弦余项公式中取 $m=1$（即 $2$ 阶展开），由 $|\cos(\theta x)|\le1$ 得 $|\sin x-x|\le |x|^3/3!$。为使右端不超过 $10^{-3}$，只需取
\[
 |x|\le\sqrt[3]{6\cdot10^{-3}}=0.181712\ldots\quad\text{（弧度）}.
\]

用三次式 $x-x^3/6$ 近似时，因四次项系数同样为零，取 $m=2$（即 $4$ 阶展开），有 $|\sin x-(x-x^3/6)|\le |x|^5/5!$。要使误差不超过 $10^{-3}$，只需满足
\[
 |x|\le\sqrt[5]{120\cdot10^{-3}}=0.654389\ldots\quad\text{（弧度）}.
\]

上述两个上界换算为角度分别约为 $10.4113^\circ$ 与 $37.4938^\circ$；若取较为简短且仍严格保证精度的弧度界，可分别写为 $|x|\le0.1817$ 与 $|x|\le0.6543$。由此可见，在相同的误差限度下，三次近似多项式显著拓宽了自变量的可用范围。
\end{proof}

\end{document}
